% -*- mode : latex; coding: utf-8 -*-
\chapter{System Design and Implementation}
\label{sec:system-design}

This chapter describes the system design and implementation of the evaluation platform used to analyze MPU performance in automotive ECU environments. The platform consists of STM32F446 hardware, custom Rust-based RTOS implementation, MPU configuration management, and performance measurement infrastructure designed for cycle-accurate analysis.

\section{Hardware Platform: STM32F446}

The STM32F446RE microcontroller serves as the evaluation platform, representing typical automotive-grade Cortex-M4 implementations. This device provides representative performance characteristics and MPU functionality found in production automotive ECU systems.

\subsection{Core Specifications}

The STM32F446RE provides the following specifications relevant to automotive ECU applications:

\begin{itemize}
    \item ARM Cortex-M4F processor with single-precision FPU
    \item 180MHz maximum operating frequency providing 225 DMIPS performance
    \item 512KB Flash memory for application code and configuration data
    \item 128KB SRAM for runtime data and task stacks
    \item 8-region Memory Protection Unit with sub-region disable capability
\end{itemize}

\subsection{Automotive-Relevant Features}

The STM32F446 includes features specifically designed for automotive environments:

\textbf{Operating Conditions}: Extended temperature range (-40°C to +105°C) and high electromagnetic compatibility (EMC) robustness ensure reliable operation in harsh automotive environments.

\textbf{Communication Interfaces}: CAN bus interface enables integration with automotive networks, while multiple UART, SPI, and I2C interfaces support sensor and actuator communication typical in automotive applications.

\textbf{Power Management}: Advanced power management capabilities including multiple low-power modes support automotive power budget requirements and enable efficient operation during vehicle standby states.

\subsection{Development Platform Configuration}

The evaluation uses an STM32F446RE Nucleo development board configured for representative automotive operation:

\textbf{Clock Configuration}: The system operates at maximum 180MHz frequency using an external high-speed crystal oscillator, ensuring timing accuracy required for cycle-precise measurements.

\textbf{Memory Layout}: The memory configuration follows automotive ECU conventions with clear separation between program code, runtime data, and peripheral memory regions to enable meaningful MPU protection analysis.

\section{Rust-based RTOS Implementation}

This research implements a custom RTOS in Rust, a systems programming language offering memory safety guarantees that complement hardware-based protection. The RTOS design follows automotive best practices while enabling detailed performance analysis.

\subsection{Language Choice Rationale}

Rust provides several advantages for automotive ECU development that justify its selection for this research:

\textbf{Memory Safety}: Compile-time memory safety verification eliminates entire classes of vulnerabilities, complementing hardware MPU protection with software-level guarantees.

\textbf{Zero-Cost Abstractions}: Rust's abstraction mechanisms compile to efficient machine code without runtime overhead, essential for meeting automotive performance requirements.

\textbf{Deterministic Behavior}: Absence of garbage collection and predictable memory allocation patterns support real-time automotive applications requiring deterministic timing behavior.

\subsection{RTOS Architecture}

The RTOS implements essential functionality required for automotive ECU operation:

\begin{algorithm}[h]
\SetAlgoLined
\KwResult{Task scheduling and MPU management}
\SetKwFunction{FSchedule}{schedule}
\SetKwFunction{FContextSwitch}{context\_switch}
\SetKwFunction{FConfigMPU}{configure\_mpu}

\SetKwProg{Fn}{Function}{:}{}
\Fn{\FSchedule{current\_task}}{
    next\_task $\leftarrow$ priority\_queue.pop()\;
    \If{next\_task.mpu\_config $\neq$ current\_task.mpu\_config}{
        \FConfigMPU{next\_task.mpu\_config}\;
    }
    \FContextSwitch{current\_task, next\_task}\;
}
\caption{Task Scheduling with MPU Management}
\end{algorithm}

\textbf{Task Management}: The task management system implements priority-based preemptive scheduling with round-robin scheduling within priority levels. Each task includes metadata for MPU configuration, enabling automatic memory protection setup during context switches.

\textbf{Memory Management}: Static memory allocation eliminates heap fragmentation and provides deterministic memory access patterns. Task stacks are allocated at compile time with MPU guard regions to detect stack overflow conditions.

\textbf{Synchronization Primitives}: The RTOS provides mutexes, semaphores, and message queues with priority inheritance support to prevent priority inversion in safety-critical automotive applications.

\subsection{Hardware Abstraction Layer}

The hardware abstraction layer (HAL) provides safe interfaces to STM32F446 peripherals while enabling performance measurement:

\begin{lstlisting}[language=Rust, caption=MPU Configuration Interface]
pub struct MpuManager {
    regions: [Option<MpuRegion>; 8],
    active_config: MpuConfiguration,
}

impl MpuManager {
    pub unsafe fn configure_region(
        &mut self,
        id: RegionId,
        region: MpuRegion
    ) -> Result<(), MpuError> {
        // Configure MPU region with cycle counting
        let start_cycles = dwt::get_cycles();
        self.write_mpu_registers(id, &region);
        let end_cycles = dwt::get_cycles();

        // Record configuration overhead
        self.record_config_overhead(end_cycles - start_cycles);
        Ok(())
    }
}
\end{lstlisting}

\section{MPU Configuration and Management}

The MPU implementation provides comprehensive memory protection while enabling detailed performance analysis of protection overhead.

\subsection{Memory Region Strategy}

The MPU configuration strategy balances security requirements with performance considerations:

\textbf{Region 0 - Flash Memory}:
\begin{itemize}
    \item Address Range: 0x08000000 - 0x0807FFFF (512KB)
    \item Attributes: Cacheable, Execute/Read permissions
    \item Purpose: Program code protection with execute permission enforcement
\end{itemize}

\textbf{Region 1 - System RAM}:
\begin{itemize}
    \item Address Range: 0x20000000 - 0x2001FFFF (128KB)
    \item Attributes: Cacheable, Read/Write permissions (no execute)
    \item Purpose: Data memory with execute prevention for security
\end{itemize}

\textbf{Region 2 - Peripheral Memory}:
\begin{itemize}
    \item Address Range: 0x40000000 - 0x5FFFFFFF (512MB)
    \item Attributes: Device memory, Read/Write permissions
    \item Purpose: Hardware register access with appropriate ordering guarantees
\end{itemize}

\textbf{Regions 3-7 - Task-Specific Protection}:
Dynamic regions configured per task for stack isolation and data protection, enabling fine-grained memory isolation between automotive software components.

\subsection{Dynamic MPU Reconfiguration}

Context switches require efficient MPU reconfiguration to maintain protection while minimizing overhead:

\begin{lstlisting}[language=Rust, caption=Context Switch MPU Update]
unsafe fn context_switch_with_mpu(
    from_task: &mut TaskContext,
    to_task: &TaskContext
) {
    // Measure context switch timing
    let start_cycles = dwt::get_cycles();

    // Save current task state
    save_task_registers(from_task);

    // Update MPU if configuration differs
    if from_task.mpu_config != to_task.mpu_config {
        mpu_disable();
        configure_task_regions(&to_task.mpu_config);
        mpu_enable();
    }

    // Load new task state
    load_task_registers(to_task);

    let end_cycles = dwt::get_cycles();
    record_context_switch_time(end_cycles - start_cycles);
}
\end{lstlisting}

\subsection{Optimization Strategies}

Several optimization techniques minimize MPU overhead while maintaining security properties:

\textbf{Lazy Reconfiguration}: MPU regions are updated only when task memory requirements differ, avoiding unnecessary reconfiguration overhead for tasks with identical memory layouts.

\textbf{Region Caching}: Frequently used MPU configurations are cached to reduce setup time during context switches, particularly beneficial for periodic automotive control tasks.

\textbf{Batch Configuration}: Multiple MPU regions are configured simultaneously during single MPU disable periods, minimizing the time window where memory protection is inactive.

\section{Performance Measurement Infrastructure}

Accurate performance measurement requires specialized infrastructure that provides cycle-accurate timing without introducing measurement artifacts.

\subsection{DWT-based Cycle Counting}

The measurement infrastructure utilizes the ARM Cortex-M4 DWT unit for precise timing:

\begin{lstlisting}[language=Rust, caption=DWT Cycle Counter Implementation]
mod dwt {
    const DWT_CTRL: *mut u32 = 0xE0001000 as *mut u32;
    const DWT_CYCCNT: *mut u32 = 0xE0001004 as *mut u32;

    pub unsafe fn init() {
        // Enable DWT cycle counter
        let mut ctrl = ptr::read_volatile(DWT_CTRL);
        ctrl |= 0x1; // CYCCNTENA bit
        ptr::write_volatile(DWT_CTRL, ctrl);

        // Reset counter to zero
        ptr::write_volatile(DWT_CYCCNT, 0);
    }

    #[inline(always)]
    pub fn get_cycles() -> u32 {
        unsafe { ptr::read_volatile(DWT_CYCCNT) }
    }

    pub fn cycles_to_microseconds(cycles: u32) -> f32 {
        cycles as f32 / 180.0 // 180MHz operation
    }
}
\end{lstlisting}

\subsection{Statistical Analysis Framework}

The measurement framework includes comprehensive statistical analysis capabilities:

\begin{lstlisting}[language=Rust, caption=Benchmark Statistics]
pub struct BenchmarkStats {
    pub samples: Vec<u32>,
    pub min_cycles: u32,
    pub max_cycles: u32,
    pub mean_cycles: f64,
    pub std_dev: f64,
    pub percentile_95: u32,
    pub percentile_99: u32,
}

impl BenchmarkStats {
    pub fn from_measurements(measurements: &[u32]) -> Self {
        let mut sorted = measurements.to_vec();
        sorted.sort_unstable();

        let mean = measurements.iter().sum::<u32>() as f64
                  / measurements.len() as f64;

        let variance = measurements.iter()
            .map(|&x| (x as f64 - mean).powi(2))
            .sum::<f64>() / measurements.len() as f64;

        Self {
            samples: measurements.to_vec(),
            min_cycles: sorted[0],
            max_cycles: sorted[sorted.len() - 1],
            mean_cycles: mean,
            std_dev: variance.sqrt(),
            percentile_95: sorted[sorted.len() * 95 / 100],
            percentile_99: sorted[sorted.len() * 99 / 100],
        }
    }
}
\end{lstlisting}

\subsection{Data Collection and Analysis}

Performance data collection follows rigorous methodology to ensure statistical validity:

\textbf{Sample Size}: Each benchmark collects minimum 100 samples to ensure statistical significance, with larger sample sizes (1000+) for operations with high variance.

\textbf{Warm-up Procedures}: All benchmarks include warm-up phases to stabilize processor caches and pipeline state before measurement collection begins.

\textbf{Outlier Detection}: Statistical analysis includes outlier detection and filtering to identify measurement artifacts that could skew results.

\textbf{Real-time Logging}: Performance data is collected using Real-Time Transfer (RTT) for external analysis, enabling detailed offline analysis without impacting real-time measurement accuracy.

This comprehensive system design provides the foundation for accurate MPU performance analysis in automotive ECU environments, combining representative hardware with optimized software implementation and precise measurement capabilities.